\documentclass[12pt,a4paper]{article}

\usepackage[utf8]{inputenc}
\usepackage[T1]{fontenc}
\usepackage{amsmath,amssymb,amsthm}
\usepackage{mathrsfs}
\usepackage{geometry}
\usepackage{hyperref}
\usepackage{booktabs}
\usepackage{array}
\usepackage{longtable}
\usepackage{enumitem}

\geometry{margin=1in}

\theoremstyle{plain}
\newtheorem{theorem}{Theorem}[section]
\newtheorem{proposition}[theorem]{Proposition}
\newtheorem{lemma}[theorem]{Lemma}
\newtheorem{corollary}[theorem]{Corollary}

\theoremstyle{definition}
\newtheorem{definition}[theorem]{Definition}
\newtheorem{axiom}[theorem]{Axiom}
\newtheorem{example}[theorem]{Example}

\theoremstyle{remark}
\newtheorem{remark}[theorem]{Remark}

\newcommand{\vol}{\operatorname{vol}}
\newcommand{\Aut}{\operatorname{Aut}}
\newcommand{\spec}{\operatorname{spec}}
\newcommand{\Tr}{\operatorname{Tr}}
\newcommand{\real}{\mathbb{R}}
\newcommand{\hilb}{\mathcal{H}}
\newcommand{\alg}{\mathcal{A}}
\newcommand{\lag}{\mathcal{L}}

\title{Constructive Dynamical Axioms for Quantum Field Theory\\[6pt]
\large Version VII: Core Axioms, Specialization Axes, Spectral Type,\\
Scope Boundaries, and Detailed Framework Correspondences}

\author{Guillem Duran Ballester\\
Fragile Technologies S.L.\\
\texttt{guillem@fragile.tech}}

\date{April 2026}

\begin{document}
\maketitle

\begin{abstract}
The standard axiomatic frameworks for quantum field theory ---
Wightman, Osterwalder--Schrader, and Haag--Kastler / locally covariant AQFT ---
describe the output of a quantum field theory rather than the
constructive dynamical system that produces it.
The present paper proposes a framework whose primitive object is the
constructive dynamical system itself, organized into five core axioms
shared by all constructive dynamical QFTs with propagating local
degrees of freedom --- geometric locality, finite resolution, finite
propagation, local dynamics with additive decomposition, and
positivity/stability --- together with four specialization axes for
internal symmetry, statistics, interaction type, and spectral type.
The framework is designed so that the standard axiomatic frameworks
reappear as specialization branches under explicitly stated
supplementary hypotheses: in the flat stationary sector, the core
axioms together with branch data and transfer-matrix hypotheses yield
the Osterwalder--Schrader axioms OS0--OS4 and hence, by
reconstruction, the Wightman axioms; on general globally hyperbolic
manifolds, the same core-plus-branch structure induces a locally
covariant net of physical $*$-algebras. The paper isolates upstream
structural conditions and maps them in detail to the standard
frameworks, separating universal framework statements from the
additional analytic hypotheses needed for clustering,
transfer-matrix factorization, hyperbolicity, and Hadamard
regularity. Its object is the abstract constructive architecture
itself.
\end{abstract}

\tableofcontents
\newpage

%% ===================================================================
\section{Introduction}
%% ===================================================================

\subsection{Problem Statement}

The mathematical foundations of quantum field theory rest on three
standard axiomatic frameworks, each native to a particular geometric
setting.

The Wightman axioms~\cite{Wightman1956, StreaterWightman} take as
primitive operator-valued tempered distributions on Minkowski space
acting on a Hilbert space carrying a positive-energy representation of
the Poincar\'e group. They encode locality, relativistic covariance,
and the spectral condition, but they are stated in terms of the output
of the theory rather than in terms of a constructive dynamical system
that generates that output.

The Osterwalder--Schrader axioms~\cite{OS1, OS2} take as primitive
the Schwinger functions on flat Euclidean space and impose
temperedness, Euclidean covariance, reflection positivity, clustering,
and graded symmetry. They are the gateway from Euclidean construction
to Minkowskian QFT, but again they operate downstream, at the level
of correlation functions already given as output data.

The Haag--Kastler axioms and their locally covariant
generalization~\cite{Haag, BFV} take as primitive a net, or functor,
of local $*$-algebras over globally hyperbolic spacetimes. These
frameworks are manifold-native and conceptually powerful, but they
organize observables rather than specifying the constructive dynamics
that produces them.

These frameworks are complementary rather than competing. The problem
is that they are fragmented. There is no widely used upstream axiom
system whose primitive object is the constructive dynamical system
itself and which is explicitly designed to specialize into all three
standard frameworks.

\subsection{Main Idea}

The framework has two layers.

\begin{itemize}[nosep]
\item \textbf{Core axioms}: structural requirements shared by all
  constructive QFTs, independent of field content or phase.
\item \textbf{Specialization axes}: branch conditions that classify
  theory type and selected phase.
\end{itemize}

The five core axioms are:

\begin{enumerate}[nosep]
\item Geometric locality;
\item Finite resolution;
\item Finite propagation;
\item Local dynamics with additive decomposition;
\item Positivity and stability.
\end{enumerate}

The specialization data are organized along four axes:

\begin{description}[nosep]
\item[Axis A:] Internal symmetry (gauge, global, or discrete gauge);
\item[Axis B:] Statistics (bosonic, fermionic, mixed, or
  supersymmetric mixed);
\item[Axis C:] Interaction type (interacting or free);
\item[Axis D:] Spectral type (massive/gapped or critical/conformal).
\end{description}

The distinction between a gapped theory and a critical/conformal
theory is not built into the core axioms; it is selected by parameter
choice or phase selection inside the same constructive architecture.

\subsection{Nature of This Paper}

This is a framework paper. The core axioms isolate upstream
structural conditions. The specialization theorems show that, under
explicitly stated supplementary hypotheses, the standard frameworks
reappear as abstract branches. The paper does not claim that the core
axioms automatically produce those hypotheses. It sharply separates
three layers of verification:

\begin{enumerate}[nosep]
\item \textbf{Structural verification}: the core axioms and branch
  data;
\item \textbf{Specialization verification}: the flat-stationary or
  manifold-covariant branch theorem;
\item \textbf{Supplementary analytic hypotheses}: transfer-matrix
  factorization, spectral gap or mixing input, hyperbolicity for the
  time-slice property, Hadamard regularity, and related conditions.
\end{enumerate}

The paper cleanly identifies which upstream input produces which
downstream property and which steps require additional analytic
hypotheses.

\subsection{Main Theorems}

The paper is organized around three results, stated here in a form
that makes the additional hypotheses explicit.

\begin{theorem}[Flat Stationary Specialization]\label{thm:flat}
  In the flat stationary sector $(M,g) = (\real^4, \delta)$, a system
  satisfying Core Axioms~1--5 together with branch choices on Axes A,
  B, and C produces Schwinger functions satisfying:
  \begin{enumerate}[nosep,label=(\roman*)]
  \item OS0 from Core Axiom~2;
  \item OS1 from Core Axiom~4 in the flat Euclidean sector, under the
    hypothesis that the equilibrium measure is the unique
    $E(4)$-invariant measure associated with the action;
  \item OS2 from Core Axioms~4 and~5, under the transfer-matrix
    hypotheses of Proposition~\ref{prop:transfer};
  \item OS4 from the chosen branch on Axis~B.
  \end{enumerate}
  If, in addition, the translation-mixing estimate
  \eqref{eq:mixing} holds for the translation-stable observable
  algebra generating the Schwinger functions, then OS3 follows.
  Axis~D refines the expected clustering rate: gapped phases (D1)
  yield exponential decay, while critical/conformal phases (D2) yield
  polynomial decay. OS reconstruction then yields Wightman functions
  satisfying W0--W4.
\end{theorem}
\begin{proof}
  This is the summary form of Theorem~\ref{thm:flat-detail}. The
  structural derivations of OS0, OS1, OS2, and OS4 are proved there,
  and OS3 is supplied by Proposition~\ref{prop:mixing}. The final
  reconstruction step is the standard Osterwalder--Schrader
  reconstruction theorem~\cite{OS1, OS2}.
\end{proof}

\begin{theorem}[Manifold-Covariant Specialization]\label{thm:manifold}
  On a general globally hyperbolic Lorentzian manifold $(M,g)$, a
  system satisfying Core Axioms~1--5, a symmetry branch on Axis~A,
  and geometric naturality of the generating structure in the sense of
  Remark~\ref{rem:naturality} induces a locally covariant net of
  physical $*$-algebras satisfying isotony, locality, and local
  covariance. In the gauge branches (A1 or A3), the physical net is
  the invariant subalgebra; in the global-symmetry branch (A2), the
  full local algebras are physical and the global symmetry acts by
  automorphisms. If the field equations of the theory are normally
  hyperbolic or of Dirac type, then the time-slice property holds. If
  the physical states are Hadamard, then the microlocal spectrum
  condition holds.
\end{theorem}
\begin{proof}
  This is the summary form of Theorem~\ref{thm:manifold-detail}. The
  net structure, covariance morphisms, and the two additional
  conclusions are proved there from Core Axioms~1--5 together with
  Remark~\ref{rem:naturality}.
\end{proof}

\begin{theorem}[Structural Separation]\label{thm:separation}
  The five core axioms together with the four specialization axes
  determine a universal separation between structural outputs and
  supplementary analytic inputs. In particular, temperedness,
  algebraic locality, positivity, and branch-dependent grading are
  fixed at the framework level, whereas reflection positivity beyond
  hyperplane decomposition, clustering, time-slice, and microlocal
  spectrum statements require the additional hypotheses made explicit
  in Sections~\ref{sec:existing} and~\ref{sec:specialization}.
\end{theorem}
\begin{proof}
  The separation is obtained by reading the later derivations result by
  result. Core Axiom~\ref{ax:resolution} yields OS0 directly; Core
  Axiom~\ref{ax:locality} together with Axis~B determines the algebraic
  locality and grading data; Core Axiom~\ref{ax:positivity} gives the
  positive semigroup and hence the GNS positivity input; and Axis~B
  fixes graded symmetry. By contrast, Proposition~\ref{prop:transfer}
  shows that OS2 requires, in addition to the core axioms, explicit
  transfer-matrix hypotheses. Proposition~\ref{prop:mixing} shows that
  OS3 requires the translation-mixing hypothesis \eqref{eq:mixing}.
  Theorem~\ref{thm:manifold-detail} shows that local covariance uses
  the extra geometric naturality stated in
  Remark~\ref{rem:naturality}, while time-slice and microlocal
  spectrum require respectively hyperbolicity and Hadamard
  regularity. This is exactly the claimed structural/analytic
  separation.
\end{proof}

\subsection{Relation to Recent Constructive QFT}

The present framework should be positioned relative to the active
frontiers of rigorous constructive QFT. Hairer's theory of regularity
structures~\cite{Hairer2014} and the paracontrolled distributions of
Gubinelli, Imkeller, and Perkowski~\cite{GIP} have made it possible
to construct solutions to singular stochastic PDEs --- including the
dynamical $\Phi^4_3$ model --- that were previously beyond rigorous
reach. The Balaban--Dimock--Imbrie--Jaffe
program~\cite{Balaban1,Balaban2} pursues constructive Yang--Mills
theory through multiscale cluster expansions.

These programs develop concrete constructions. The present framework
is not intended to compete with them but to provide an organizational
layer above them: a common set of structural conditions and
specialization rules that can be used to compare different
constructive approaches without committing the framework itself to a
particular one.


\subsection{Roadmap}

Section~\ref{sec:existing} develops the detailed correspondence
between the present framework and the three standard axiomatic
frameworks. Section~\ref{sec:core} states the core axioms, with
particular attention to the reformulation of Core Axiom~4 and the
causal-structure role of Core Axiom~3. Section~\ref{sec:axes} states
the specialization axes. Section~\ref{sec:specialization} proves the
specialization theorems, with an explicit treatment of the upstream
dynamical chain leading to OS3. Section~\ref{sec:compatibility}
discusses compatibility, non-redundancy, and core separating examples.
Section~\ref{sec:constructive} explains in what sense the framework is
constructive and clarifies the role of renormalization and scaling.
Section~\ref{sec:susy} treats supersymmetry and the
TQFT scope boundary. Section~\ref{sec:discussion} records limitations
and future refinements. Section~\ref{sec:conclusion} concludes.


%% ===================================================================
\section{The Standard Frameworks and Their Upstream Sources}
\label{sec:existing}
%% ===================================================================

This section develops the detailed correspondence between the five
core axioms (plus specialization data) and each of the three standard
axiomatic frameworks. For each standard axiom, we identify: which core
axiom provides the primary structural input; which additional
hypotheses, if any, are needed; and where the hard analytic work
resides. This makes explicit the claim that the present framework is
genuinely upstream of the existing ones.


\subsection{Wightman Axioms}

The Wightman axioms~\cite{Wightman1956, StreaterWightman} take as
primitive a separable Hilbert space $\hilb$, a vacuum vector
$|\Omega\rangle$, a unitary strongly continuous representation of the
proper orthochronous Poincar\'e group, and a collection of
operator-valued tempered distributions obeying covariance, locality,
the spectral condition, and cyclicity of the vacuum. In the present
framework, none of these are primitive; they are all derived.

\medskip

\noindent\textbf{W0: Hilbert space, vacuum, Poincar\'e
  representation, spectral condition.} Core Axiom~5 provides a
positivity-preserving semigroup with self-adjoint generator $H \ge 0$
and an equilibrium state $\omega$. The GNS construction on $\omega$
produces $(\hilb, \pi, |\Omega\rangle)$. Uniqueness of the vacuum
follows from OS3, which requires the translation-mixing hypothesis of
Proposition~\ref{prop:mixing}. The
Poincar\'e representation arises by OS reconstruction from $E(4)$
covariance of the Schwinger functions (OS1). The spectral condition
is part of the standard Osterwalder--Schrader reconstruction output.
Core Axiom~3 is best viewed here as the upstream causal datum that is
intended to be compatible with that Lorentzian interpretation; see
Remark~\ref{rem:propagation-chain} below.

\medskip

\noindent\textbf{W1: Covariance of fields.} The field operators
transform covariantly under the Poincar\'e representation:
$U(a,\Lambda)\phi(x)U(a,\Lambda)^{-1} = S(\Lambda)\phi(\Lambda x +
a)$. This is inherited from the $E(4)$ invariance of the action
density in the flat sector (Core Axiom~4, whose local density
$\lag(\Phi, D\Phi, \delta)$ depends on the flat metric and is
therefore $E(4)$-invariant) via OS1 and analytic continuation. The
spin/statistics content comes from the branch choice on Axis~B.

\medskip

\noindent\textbf{W2: Fields as operator-valued tempered
  distributions.} This is exactly OS0, which follows from Core
Axiom~2 (finite resolution). The polynomial bounds in Schwartz
seminorms, equation~\eqref{eq:poly-bounds}, guarantee that the
Schwinger functions define tempered distributions.

\medskip

\noindent\textbf{W3: Local commutativity (microscopic causality).}
For spacelike separated points, $[\phi(x),\phi(y)]_\pm = 0$. This
combines Core Axiom~1 (locality of the algebra) with Core Axiom~3
(finite propagation defining the causal structure). After OS
reconstruction and Wick rotation, Euclidean locality becomes
Minkowski locality with the light cone defined by $c_{\mathrm{info}}$.
The $\pm$ grading is fixed by Axis~B.

\medskip

\noindent\textbf{W4: Cyclicity of the vacuum.} The set of vectors
$\{\phi(f_1)\cdots\phi(f_n)|\Omega\rangle\}$ is dense in $\hilb$.
This is automatic from the GNS construction: $|\Omega\rangle$ is
cyclic for $\pi(\alg)$ by definition.


\begin{table}[ht]
\centering
\renewcommand{\arraystretch}{1.2}
\begin{tabular}{@{}lll@{}}
\toprule
Wightman axiom & Primary source & Additional hypotheses \\
\midrule
W0 (Hilbert space, vacuum, & Core 5 (GNS) + & OS reconstruction, \\
\quad Poincar\'e, spectral cond.) & OS reconstruction & mixing / clustering input \\
W1 (covariance) & Core 4 ($E(4)$-inv.\ $\lag$) & OS reconstruction \\
 & + Axis B & \\
W2 (temperedness) & Core 2 & None \\
W3 (locality) & Core 1 + Core 3 & Axis B for grading \\
W4 (cyclicity) & Core 5 & Automatic from GNS \\
\bottomrule
\end{tabular}
\caption{Wightman axioms and their upstream sources in the present
  framework.}
\label{tab:wightman}
\end{table}


\subsection{Osterwalder--Schrader Axioms}

The OS axioms~\cite{OS1, OS2} take as primitive the Schwinger
functions $S_n(x_1,\ldots,x_n)$ on $(\real^4)^n$ and impose
temperedness (OS0), Euclidean covariance (OS1), reflection positivity
(OS2), clustering (OS3), and graded symmetry (OS4). In the present
framework the Schwinger functions are not primitive: they are produced
by the constructive dynamical system.

\medskip

\noindent\textbf{OS0 (Temperedness).} Core Axiom~2 gives polynomial
bounds on smeared observables, equation~\eqref{eq:poly-bounds}. These
directly imply that the Schwinger functions define tempered
distributions on $(\real^4)^n$. No additional hypotheses are needed.

\medskip

\noindent\textbf{OS1 (Euclidean covariance).} Core Axiom~4 provides a
local generating structure whose density, in the flat sector, depends
on fields, their derivatives, and the flat metric $\delta_{\mu\nu}$.
Since $\delta$ is $E(4)$-invariant, the Boltzmann weight
$e^{-S}$ is $E(4)$-invariant, and so is the equilibrium measure
$\pi$. Hence all expectations under $\pi$ --- the Schwinger functions
--- are $E(4)$-invariant.

\emph{Additional hypothesis}: the equilibrium measure must be uniquely
determined as an $E(4)$-invariant measure. If there are multiple
equilibrium measures (as in spontaneous symmetry breaking), one must
select an $E(4)$-invariant one.

\medskip

\noindent\textbf{OS2 (Reflection positivity).} This is the most
complex derivation. Core Axiom~4 supplies the additive decomposition
$S = S_+ + S_- + B$ across the reflection hyperplane
(Lemma~\ref{lem:hyperplane}). This lets the functional integral
factorize across the hyperplane. Core Axiom~5 gives a
positivity-preserving semigroup $e^{-tH}$ with $H \ge 0$. The
transfer-matrix identification (Proposition~\ref{prop:transfer}) then
gives $\langle \Theta F, F\rangle_E = \|e^{-t_{\min}H}F\|^2 \ge 0$.

\emph{Additional hypotheses}: the transfer-matrix hypotheses of
Proposition~\ref{prop:transfer} --- that the Euclidean measure
factorizes in the right way across the hyperplane. This is
supplementary analytic input and is where the constructive hard work
happens in concrete theories.

Core Axiom~4 (additive decomposition) is structurally indispensable
here. Without it, the transfer-matrix argument has no starting point.

\medskip

\noindent\textbf{OS3 (Clustering).} The theorem-level hypothesis
needed for OS3 in the abstract framework is the translation-mixing
estimate
\begin{equation}
  \big|\langle F \, \tau_a G\rangle_\pi
  - \langle F\rangle_\pi \langle G\rangle_\pi\big|
  \xrightarrow{|a| \to \infty} 0
\end{equation}
for the translation-stable observable algebra from which the
Schwinger functions are built. Proposition~\ref{prop:mixing} below
shows that this estimate is exactly what is needed to derive OS3.

In concrete theories one often tries to obtain~\eqref{eq:mixing} from
stronger dynamical input such as a spectral gap for an appropriate
Euclidean transfer operator or Markov generator. That route is
model-dependent and therefore lies outside the abstract theorem. The
tools used there are probabilistic --- Foster--Lyapunov drift
conditions, coupling arguments, log-Sobolev inequalities,
Bakry--\'Emery criteria, hypocoercivity --- rather than the Euclidean
field-theory tools traditionally used for clustering.

Axis~D does not affect the \emph{existence} of OS3. It refines the
\emph{rate}: under D1 (gapped), exponential clustering
$S^c \sim e^{-\Delta|a|}$; under D2 (critical), polynomial clustering
$S^c \sim |a|^{-2\Delta_\phi}$.

\medskip

\noindent\textbf{OS4 (Graded symmetry).} Purely determined by the
branch choice on Axis~B. B1 gives full symmetry, B2 gives fermionic
antisymmetry, B3 gives the mixed graded-symmetry rule, B4 preserves
B3's grading with additional boson--fermion pairing constraints. No
additional hypotheses.


\begin{table}[ht]
\centering
\renewcommand{\arraystretch}{1.2}
\begin{tabular}{@{}lll@{}}
\toprule
OS axiom & Primary source & Additional hypotheses \\
\midrule
OS0 (temperedness) & Core 2 & None \\
OS1 (Euclidean cov.) & Core 4 (flat sector) &
Unique $E(4)$-inv.\ equilibrium \\
OS2 (reflection pos.) & Core 4 + Core 5 & Transfer-matrix hypotheses \\
OS3 (clustering) & Core 5 + Prop.~\ref{prop:mixing} &
Translation-mixing estimate~\eqref{eq:mixing} \\
OS4 (graded symmetry) & Axis B & None \\
\bottomrule
\end{tabular}
\caption{OS axioms and their upstream sources.}
\label{tab:os}
\end{table}


\subsection{Haag--Kastler / Locally Covariant AQFT}

The Haag--Kastler axioms~\cite{Haag} and their locally covariant
generalization~\cite{BFV} take as primitive a net or functor assigning
local $*$-algebras to spacetime regions or spacetimes. In the present
framework this is the most natural specialization because Core Axiom~1
is already stated in Haag--Kastler language.

\medskip

\noindent\textbf{HK1 (Isotony).} $O_1 \subset O_2$ implies
$\alg(O_1) \subset \alg(O_2)$. This is literally the first clause of
Core Axiom~1. No additional hypotheses.

\medskip

\noindent\textbf{HK2 (Locality / Einstein causality).} For spacelike
separated $O_1, O_2$, the algebras commute (with appropriate grading
from Axis~B). This is the second clause of Core Axiom~1, with causal
disjointness defined by the Lorentzian metric in the manifold branch.
Core Axiom~3 ensures that the dynamical notion of causality is
compatible with the metric notion.

\medskip

\noindent\textbf{HK3 (Covariance / local covariance).} In the
Brunetti--Fredenhagen--Verch formulation, the primitive is a functor
from the category of globally hyperbolic spacetimes to $*$-algebras
with injective $*$-homomorphisms. This comes from Core Axiom~4
together with the geometric naturality stated in
Remark~\ref{rem:naturality}:
because $\lag$ depends locally on $\Phi$, $D\Phi$, and $g$, an
isometric embedding $\psi: M_1 \to M_2$ induces a pullback of field
configurations and hence a pushforward of observables, giving the
functorial structure.

\medskip

\noindent\textbf{HK4 (Time-slice property).} If $O$ contains a
Cauchy surface for a larger region $O'$, then
$\alg(O) = \alg(O')$. This is \emph{not} automatic from the core
axioms. It requires well-posed hyperbolic evolution: the field
equations must be normally hyperbolic or Dirac-type.

Core Axiom~3 is necessary but not sufficient: finite propagation
ensures that information \emph{stays within} the domain of
dependence, but time-slice additionally requires that it \emph{fills}
the domain of dependence. Finite propagation prevents leaking out;
hyperbolicity ensures filling in.

\medskip

\noindent\textbf{HK5 (Spectrum condition / positivity).} This comes
from Core Axiom~5 ($H \ge 0$) together with Core Axiom~3 (finite
propagation $\to$ bounded dispersion $\to$ spectral support in forward
cone).

\medskip

\noindent\textbf{Microlocal spectrum condition.} On curved
spacetimes, the spectrum condition is replaced by the Hadamard
condition on the two-point function. This requires the additional
hypothesis that the physical states are Hadamard --- the
curved-spacetime analog of the transfer-matrix hypothesis.

\medskip

\noindent\textbf{Physical algebra selection.} Which algebra counts as
physical depends on Axis~A: A1 or A3 (gauge) gives the invariant
subalgebra $\alg^G_M(O)$ or $\alg^{G_{\mathrm{disc}}}_M(O)$; A2
(ungauged) gives the full algebra $\alg_M(O)$ with global symmetries
acting as automorphisms.

\begin{table}[ht]
\centering
\renewcommand{\arraystretch}{1.2}
\begin{tabular}{@{}lll@{}}
\toprule
HK/LC property & Primary source & Additional hypotheses \\
\midrule
Isotony & Core 1 (clause i) & None \\
Locality & Core 1 (clause ii) + Core 3 & Axis B for grading \\
Covariance / functoriality & Core 4 (local geom.\ $\lag$) &
Geometric naturality \\
Time-slice & Core 3 (necessary) & Hyperbolicity of field eqs.\ \\
Spectrum condition & Core 5 + Core 3 & None (flat sector) \\
Microlocal spectrum cond. & Core 5 + Core 3 & Hadamard state hyp. \\
Physical algebra & Axis A & None \\
\bottomrule
\end{tabular}
\caption{Haag--Kastler / locally covariant AQFT properties and their
  upstream sources.}
\label{tab:hk}
\end{table}


\subsection{Summary: What the Core Axioms Provide}

The five core axioms are upstream of all three frameworks
simultaneously:

\begin{itemize}[nosep]
\item Core~1 $\to$ algebraic infrastructure (nets, isotony, locality)
  $\to$ shared by HK, implied by OS;
\item Core~2 $\to$ regularity (temperedness, polynomial bounds) $\to$
  OS0, W2;
\item Core~3 $\to$ causal structure $\to$ supplies the intended
  upstream causal input for the Euclidean/Lorentzian correspondence;
\item Core~4 $\to$ dynamical generation (additive decomposition) $\to$
  hyperplane splitting for OS2; together with
  Remark~\ref{rem:naturality} it gives functoriality for HK/LC;
\item Core~5 $\to$ positivity (GNS, vacuum, bounded-below energy)
  $\to$ W0, OS2, HK spectrum condition.
\end{itemize}

The additional hypotheses are cleanly separated:

\begin{itemize}[nosep]
\item Transfer-matrix hypotheses (supplementary analytic input needed
  for OS2);
\item Translation-mixing hypothesis \eqref{eq:mixing} (needed for
  OS3);
\item Geometric naturality of the generating structure (needed for
  local covariance);
\item Hyperbolicity of field equations (needed for time-slice);
\item Hadamard state condition (needed for microlocal spectrum
  condition);
\item $E(4)$-invariant equilibrium selection (needed for OS1).
\end{itemize}

This is the organizational value of the framework: in the existing
literature, these ingredients are mixed together. The Wightman axioms
assume the Hilbert space and spectral condition. The OS axioms assume
temperedness and reflection positivity. The present framework isolates
which constructive input produces which output property and clearly
marks where the hard analytic work still lives.


\begin{table}[ht]
\centering
\renewcommand{\arraystretch}{1.2}
\begin{tabular}{@{}p{3.5cm}p{3.5cm}p{3.5cm}p{3.5cm}@{}}
\toprule
Framework & Native setting & Primitive objects & Status here \\
\midrule
Wightman & Minkowski $\real^{1,3}$ &
Operator-valued fields, Hilbert space, Poincar\'e representation &
Derived (flat sector) \\
Osterwalder--Schrader & Euclidean $\real^4$ &
Schwinger functions, reflection positivity &
Derived (flat sector) \\
Haag--Kastler / LC-AQFT & General spacetime &
Nets / functors of local $*$-algebras &
Derived (manifold sector) \\
Present framework & Manifold-native &
Core structural properties + branch data &
Primitive \\
\bottomrule
\end{tabular}
\caption{The four axiomatic layers and their status in the present
  framework.}
\label{tab:layers}
\end{table}


%% ===================================================================
\section{The Five Core Axioms}
\label{sec:core}
%% ===================================================================

\subsection{Geometric Setting}

\begin{remark}[Two geometric branches]\label{rem:geom}
  The axioms are stated on an oriented manifold $(M,g)$ with metric of
  appropriate signature. In the Euclidean branch used for Schwinger
  functions, $g$ is Riemannian and the flat specialization is
  $(\real^4, \delta)$. In the Lorentzian branch used for the
  Haag--Kastler / LC-AQFT specialization, $g$ is Lorentzian and
  globally hyperbolic.

  In the Lorentzian branch, causal disjointness is the usual metric
  notion. In the Euclidean branch there is no metric light cone;
  causal disjointness is defined operationally by Core Axiom~3
  relative to the constructive dynamics. The locality clause of Core
  Axiom~1 is therefore logically downstream of Core Axiom~3 in the
  Euclidean branch.
\end{remark}


\subsection{Why a Hierarchical Axiom System}

A scalar theory such as $\lambda\phi^4$ and a gauge theory such as
Yang--Mills should agree on locality, finite resolution, finite
propagation, local dynamics, and positivity, but differ on the
symmetry axis. A free scalar and a free Dirac field should agree on
the core axioms and differ only on the statistics axis. A massive
phase and its tuned critical counterpart should agree on the same
core axioms and branch choices on A, B, C, but differ on the
infrared spectral axis~D.

The logical structure is:
$\text{core} \longrightarrow \text{specialization by branch choice}.$


\subsection{Core Axiom 1: Geometric Locality}

\begin{axiom}[Geometric Locality]\label{ax:locality}
  For each oriented manifold $(M,g)$ with metric of appropriate
  signature, there exists a net of local observable $*$-algebras
  $\alg_M(O)$ indexed by open regions $O \subset M$, satisfying:
  \begin{enumerate}[nosep,label=(\roman*)]
  \item \textbf{Isotony}: if $O_1 \subset O_2$, then
    $\alg_M(O_1) \subset \alg_M(O_2)$;
  \item \textbf{Locality}: if $O_1$ and $O_2$ are causally disjoint,
    then $[\alg_M(O_1), \alg_M(O_2)]_\pm = 0$.
  \end{enumerate}
  Each $\alg_M(O)$ is a unital $*$-algebra. In the Lorentzian branch
  it admits a $C^*$-completion. In the Euclidean branch it is
  represented as a unital $*$-subalgebra of $L^\infty(\pi)$ for the
  equilibrium measure~$\pi$.
\end{axiom}


\subsection{Core Axiom 2: Finite Resolution}

\begin{axiom}[Finite Resolution]\label{ax:resolution}
  There exists a strictly positive resolution scale $\ell_L > 0$ such
  that no operationally distinguishable states resolve below
  $\ell_L$. Smeared observables
  \begin{equation}\label{eq:smear}
    O_f = \int f(x)\, O(x)\, d\vol_g(x),
    \qquad f \in C^\infty_c(M),
  \end{equation}
  are well-defined. In the flat Euclidean sector
  $(M,g) = (\real^4, \delta)$ they satisfy polynomial bounds
  \begin{equation}\label{eq:poly-bounds}
    |\langle O_{f_1} \cdots O_{f_n}\rangle|
    \le C^n \prod_{k=1}^n \|f_k\|_{n_k},
  \end{equation}
  uniform in $\ell_L$, where $\|\cdot\|_n$ denotes a Schwartz
  seminorm.
\end{axiom}

The axiom is structural. It says the theory is formulated at finite
operational resolution. It is deliberately silent on what scaling
limit, if any, is taken along a scale family.


\subsection{Core Axiom 3: Finite Propagation}

\begin{axiom}[Finite Propagation]\label{ax:propagation}
  There exists a maximum information speed $c_{\mathrm{info}} > 0$
  such that causal influence is restricted to the domain of dependence
  determined by $c_{\mathrm{info}}$. For any perturbation localized in
  a region $O$ at time $t_0$, the resulting change in observables at
  time $t > t_0$ is supported inside the
  $c_{\mathrm{info}}(t-t_0)$-neighborhood of~$O$.
\end{axiom}

\begin{remark}[Dynamics versus equilibrium correlations]\label{rem:dynamics-vs-eq}
  This axiom concerns the \emph{dynamics} of the constructive system,
  not the \emph{equilibrium correlations} (Schwinger functions). A
  nearest-neighbor Ising model has finite-range dynamics, but the
  equilibrium two-point function $\langle \sigma_x \sigma_y\rangle$
  extends to all distances via exponential decay. No contradiction:
  the equilibrium measure integrates over all histories of the
  dynamics, including arbitrarily long chains of nearest-neighbor
  interactions.
\end{remark}

\begin{remark}[Interpretive role of finite propagation]
  \label{rem:propagation-chain}
  Finite propagation is the framework's intended upstream causal datum.
  In the Euclidean branch, where the metric has no light cone, it gives
  operational content to causal disjointness and therefore to the
  locality clause of Core Axiom~1. In concrete constructions one
  expects this dynamical notion of finite information speed to be
  compatible with the Lorentzian causal structure appearing after
  Osterwalder--Schrader reconstruction.

  This remark is interpretive rather than theorem-level. The actual
  derivation of the Wightman spectral condition in the paper comes from
  the OS reconstruction theorem once OS0--OS4 have been established,
  not from Core Axiom~3 alone.
\end{remark}

\begin{remark}[Causality is imposed at the dynamical level]
  Finite propagation is a property of the generating dynamics
  themselves. Causality is not an abstract requirement imposed after
  the fact on correlation functions; it is part of the framework at
  the level of information propagation.
\end{remark}


\subsection{Core Axiom 4: Local Dynamics with Additive Decomposition}

\begin{axiom}[Local Dynamics with Additive Decomposition]\label{ax:action}
  The dynamics are generated by a local functional that decomposes
  additively over regions. Concretely, there exists a functional $S$
  --- or, more generally, a local generating structure for the
  Euclidean measure or evolution law --- such that for
  disjoint subregions $O_1, O_2 \subset M$,
  \begin{equation}\label{eq:additive}
    S[\Phi|_{O_1 \cup O_2}]
    = S[\Phi|_{O_1}] + S[\Phi|_{O_2}] + B[\Phi|_\partial],
  \end{equation}
  where $B$ depends only on interface data. The local algebra
  $\alg_M(O)$ is generated by observables supported in~$O$.
\end{axiom}

\begin{remark}[The standard Lagrangian case]\label{rem:lagrangian}
  In most constructive QFTs --- scalar, gauge, fermionic,
  supersymmetric --- the local generating structure is a classical
  action functional
  \begin{equation}\label{eq:classical-action}
    S = \int_M \lag(\Phi, D\Phi, g)\, d\vol_g
  \end{equation}
  built from a Lagrangian density. For these, the axiom reduces to the
  standard statement that the action is the spacetime integral of a
  local density.
\end{remark}

\begin{remark}[General local generators]\label{rem:stochastic-action}
  The axiom is stated more broadly than the classical Lagrangian case.
  It allows any local generating structure whose restriction to
  separated regions obeys the additive decomposition
  \eqref{eq:additive}. The point of the axiom is locality of
  generation, not commitment to a single formalism.
\end{remark}

\begin{remark}[Geometric naturality for the manifold branch]
  \label{rem:naturality}
  The locally covariant specialization requires, in addition to
  additivity, naturality under orientation-preserving isometric
  embeddings. Concretely, if $\psi: M_1 \hookrightarrow M_2$ is such
  an embedding, then pullback of fields by~$\psi$ must intertwine the
  generating structures, and the induced pushforward of observables
  must define injective $*$-homomorphisms
  \[
    \alpha_\psi: \alg_{M_1}(O) \to \alg_{M_2}(\psi(O))
  \]
  satisfying $\alpha_{\psi_2 \circ \psi_1}
  = \alpha_{\psi_2}\circ \alpha_{\psi_1}$ and
  $\alpha_{\mathrm{id}_M} = \mathrm{id}_{\alg_M}$. This is the
  additional geometric input used in the manifold-covariant
  specialization theorem.
\end{remark}

\begin{remark}[Scope exclusions]\label{rem:action-scope}
  Theories without any local generating structure --- self-dual field
  theories lacking even an auxiliary action formulation (such as the
  chiral 2-form on a 6-manifold), or theories defined purely through
  algebraic or functorial data without a constructive Euclidean
  formulation --- are outside the framework's scope. This is a
  deliberate restriction: the transfer-matrix route to OS2 requires
  additive decomposition as a structural input, and any theory that
  lacks such decomposition cannot enter the flat-sector specialization
  through the present architecture. This parallels the TQFT scope
  boundary of Section~\ref{sec:tqft}: the exclusion is stated openly.

  Theories that \emph{have} a classical action but face difficulties
  in defining the path integral measure --- such as chiral gauge
  theories with anomalies --- are not excluded by Core Axiom~4. Their
  difficulties lie elsewhere: in the well-definedness of the
  gauge-invariant integration and hence in the verification of
  Axis~A1, not in the locality of the generating structure.
\end{remark}


\subsection{Core Axiom 5: Positivity and Stability}

\begin{axiom}[Positivity and Stability]\label{ax:positivity}
  There exists a positivity-preserving semigroup $e^{-tH}$ on the
  algebra designated as physical by the chosen symmetry branch on
  Axis~A, with self-adjoint generator $H \ge 0$. Together with a
  chosen equilibrium state~$\omega$, this yields a GNS Hilbert-space
  representation $(\hilb, \pi, |\Omega\rangle)$ with
  $H|\Omega\rangle = 0$.
\end{axiom}

The core requires bounded-below energy and positivity of Euclidean
time evolution. It does not require a spectral gap. Whether the
selected phase is gapped or critical belongs to Axis~D.


\subsection{Finite Resolution and Scaling}

Core Axiom~2 remains a core axiom rather than a branch condition. The
framework is formulated at strictly positive operative resolution
$\ell_L > 0$, and questions about continuum limits, critical limits,
or scale families are separate analytic questions layered on top of
that structural starting point.

If one studies a scale-labelled family of theories
$\{\mathcal{T}_\lambda\}_{\lambda \in \Lambda}$, each member is
evaluated against the same core axioms and branch logic. The
renormalization group then compares different scales inside such a
family without changing the logical status of the axioms themselves.

Thus finite resolution is not a disposable regulator inside the
framework. It is the structural datum from which one may later ask
whether any additional scaling theorem exists.


%% ===================================================================
\section{The Specialization Axes}
\label{sec:axes}
%% ===================================================================

A theory satisfying the core axioms becomes a concrete branch of
constructive QFT only after branch choices are specified. Axes A--C
classify field content and interaction structure. Axis~D classifies
the infrared phase.


\subsection{Axis A: Internal Symmetry}

\begin{axiom}[A1: Gauge-Invariant Physical Algebra]\label{ax:A1}
  There exists a compact gauge group $G$ acting locally on fields. The
  physical observable algebra is the gauge-invariant subalgebra
  $\alg^{\mathrm{phys}}_M(O) = \alg_M(O)^G$.
\end{axiom}

\emph{Scope.} Yang--Mills with $G = SU(N)$, electrodynamics,
electroweak theory, Standard Model gauge sector.

\begin{axiom}[A2: Global Symmetry / Ungauged Theory]\label{ax:A2}
  The theory carries a global symmetry group $G_{\mathrm{glob}}$
  (possibly trivial) but no local gauge redundancy:
  $\alg^{\mathrm{phys}}_M(O) = \alg_M(O)$.
\end{axiom}

\emph{Scope.} $\lambda\phi^4$, Ising field theory, $O(N)$ models,
sigma models, free theories without gauge fields.

\begin{axiom}[A3: Discrete Gauge Theory]\label{ax:A3}
  There exists a finite gauge group $G_{\mathrm{disc}}$ acting locally
  on fields:
  $\alg^{\mathrm{phys}}_M(O) = \alg_M(O)^{G_{\mathrm{disc}}}$.
\end{axiom}

\emph{Scope.} $\mathbb{Z}_N$ gauge theories, Dijkgraaf--Witten type
theories.


\subsection{Axis B: Statistics}

\begin{axiom}[B1: Bosonic Statistics]\label{ax:B1}
  The field algebra carries a trivial $\mathbb{Z}_2$ grading.
  Schwinger functions are symmetric under permutation.
\end{axiom}

\begin{axiom}[B2: Fermionic Sector]\label{ax:B2}
  There exists a spin structure and a $\mathbb{Z}_2$-graded field
  algebra with commutation and anticommutation at spacelike separation
  according to grading. The action contains a Dirac-type kinetic term.
\end{axiom}

\begin{axiom}[B3: Mixed Bose--Fermi Statistics]\label{ax:B3}
  The field algebra decomposes as a graded product
  $\alg \sim \alg_{\mathrm{bos}} \otimes \widehat{\alg}_{\mathrm{ferm}}$
  with combined $\mathbb{Z}_2$ grading and Yukawa-type couplings.
\end{axiom}

\begin{axiom}[B4: Supersymmetric Mixed Sector]\label{ax:B4}
  The theory satisfies Branch~B3 and, in addition, carries an odd
  supercharge $Q$ such that
  \begin{equation}\label{eq:susy}
    H = Q^\dagger Q
  \end{equation}
  on the physical Hilbert space.
\end{axiom}

\begin{proposition}\label{prop:susy-positivity}
  Under Branch~B4, the positivity $H \ge 0$ of Core Axiom~5 holds
  automatically on the common invariant core of~\eqref{eq:susy}.
\end{proposition}
\begin{proof}
  Let $\psi$ lie in the common invariant core of $Q$ and $Q^\dagger$.
  Using~\eqref{eq:susy},
  \[
    \langle \psi, H\psi\rangle
    = \langle \psi, Q^\dagger Q\psi\rangle
    = \langle Q\psi, Q\psi\rangle
    = \|Q\psi\|^2 \ge 0.
  \]
  Hence the quadratic form of~$H$ is nonnegative on the stated core,
  which is exactly the positivity required in Core
  Axiom~\ref{ax:positivity}.
\end{proof}

\begin{remark}
  Supersymmetry is not a new axis independent of statistics; it
  presupposes coexistence of bosonic and fermionic sectors and then
  adds an algebraic relation linking them. It is therefore a
  refinement of B3 rather than a fifth axis.
\end{remark}

\begin{remark}[Protected vacuum energy and the Witten index]
  Under Branch~B4, the quantity
  $\Tr(-1)^F e^{-\beta H}$ is stable under continuous
  supersymmetry-preserving deformations~\cite{Witten}. A nonzero index
  implies unbroken supersymmetry and exactly zero vacuum energy.
\end{remark}


\subsection{Axis C: Interaction Type}

\begin{axiom}[C1: Interacting Theory]\label{ax:C1}
  At least one gauge-invariant coupling constant, self-coupling, or
  curvature observable is nonzero, or equivalently there exists
  $n \ge 4$ such that $S^c_n \not\equiv 0$.
\end{axiom}

\begin{axiom}[C2: Free Theory]\label{ax:C2}
  The action is quadratic:
  $S = \int \frac{1}{2}\Phi \cdot D \cdot \Phi\, d\vol_g$. All
  connected $n$-point functions vanish for $n \ge 3$.
\end{axiom}


\subsection{Axis D: Spectral Type}

\begin{axiom}[D1: Massive / Gapped Phase]\label{ax:D1}
  In the GNS representation, $\spec(H) \cap (0, \Delta) = \emptyset$
  for some $\Delta > 0$. When the corresponding framework hypotheses
  furnish the
  corresponding clustering estimate, connected correlations decay
  exponentially.
\end{axiom}

\begin{axiom}[D2: Critical / Conformal Phase]\label{ax:D2}
  The vacuum sector is gapless:
  $\inf \spec H|_{\hilb \ominus \mathbb{C}|\Omega\rangle} = 0$. The
  long-distance sector is scale-invariant, and connected correlations
  decay polynomially.
\end{axiom}

\begin{remark}
  A fixed positive $\ell_L$ breaks exact microscopic dilation
  invariance. What becomes scale-invariant at criticality is the
  long-distance sector, not the microscopic theory. Core Axiom~2
  remains intact.
\end{remark}

\begin{remark}[Scope of Axis~D]
  The present framework isolates the two clean spectral regimes most
  relevant to the constructive program. Gapless but nonconformal
  theories (e.g., theories with massless gauge bosons) motivate a
  future Branch~D3.
\end{remark}


\subsection{Dependency Diagram and Classification Table}

\begin{table}[ht]
\centering
\renewcommand{\arraystretch}{1.2}
\begin{tabular}{@{}lcccc@{}}
\toprule
Theory / phase & A & B & C & D \\
\midrule
$\lambda\phi^4$ massive phase & A2 & B1 & C1 & D1 \\
Free massive scalar & A2 & B1 & C2 & D1 \\
Pure Yang--Mills with mass gap & A1 & B1 & C1 & D1 \\
Massive free Dirac field & A2 & B2 & C2 & D1 \\
Wess--Zumino (unbroken massive) & A2 & B4 & C1 & D1 \\
$O(N)$ sigma model massive & A2 & B1 & C1 & D1 \\
$\mathbb{Z}_2$ gauge theory & A3 & B1 & C1 & D1 \\
Critical Ising / scalar fixed point & A2 & B1 & C1 & D2 \\
Critical $O(N)$ fixed point & A2 & B1 & C1 & D2 \\
\bottomrule
\end{tabular}
\caption{Representative theories and phases in the framework.}
\label{tab:theories}
\end{table}


%% ===================================================================
\section{Specialization Theorems}
\label{sec:specialization}
%% ===================================================================

\subsection{Flat Stationary Sector}

\begin{definition}\label{def:flat}
  The flat stationary sector is the Euclidean specialization
  $(M,g) = (\real^4, \delta)$ with flat Euclidean metric,
  time-translation invariant dynamics, and Euclidean time reflection
  $\Theta: (x_0, \mathbf{x}) \mapsto (-x_0, \mathbf{x})$. The
  positive-time physical algebra $\alg^{\mathrm{phys}}_+$ consists of
  physical observables supported in $\{x_0 > 0\}$.
\end{definition}


\subsection{Hyperplane Decomposition and Transfer Matrix}

\begin{lemma}[Hyperplane decomposition]\label{lem:hyperplane}
  In the flat Euclidean sector, the local generating functional of
  Core Axiom~\ref{ax:action} decomposes relative to the reflection
  hyperplane $\{x_0 = 0\}$:
  \begin{equation}\label{eq:hyperplane}
    S_E[\Phi] = S_E[\Phi_+] + S_E[\Phi_-] + B[\Phi_0],
  \end{equation}
  where $\Phi_\pm$ denotes the restriction to $\{\pm x_0 > 0\}$ and
  $B$ depends only on boundary data at $x_0 = 0$.
\end{lemma}
\begin{proof}
  Let $O_+^\varepsilon = \{x_0 > \varepsilon\}$ and
  $O_-^\varepsilon = \{x_0 < -\varepsilon\}$ for $\varepsilon > 0$.
  These are disjoint open regions, so Core Axiom~\ref{ax:action}
  gives
  \[
    S_E[\Phi|_{O_+^\varepsilon \cup O_-^\varepsilon}]
    = S_E[\Phi|_{O_+^\varepsilon}]
    + S_E[\Phi|_{O_-^\varepsilon}]
    + B_\varepsilon[\Phi|_{\partial O_+^\varepsilon \cup \partial O_-^\varepsilon}].
  \]
  Because the generating structure is local and depends on finitely
  many derivatives, the only couplings omitted from the two bulk terms
  are those crossing the slabs $|x_0| \le \varepsilon$. As
  $\varepsilon \downarrow 0$, these couplings collapse onto the
  interface $x_0 = 0$ and therefore depend only on the boundary jet
  $\Phi_0$ there. Absorbing the limit into a single interface
  functional $B[\Phi_0]$ yields
  \[
    S_E[\Phi] = S_E[\Phi_+] + S_E[\Phi_-] + B[\Phi_0],
  \]
  which is exactly~\eqref{eq:hyperplane}.
\end{proof}

\begin{proposition}[Transfer-matrix factorization]\label{prop:transfer}
  Assume:
  \begin{enumerate}[nosep,label=(\roman*)]
  \item the generating functional satisfies the hyperplane
    decomposition of Lemma~\ref{lem:hyperplane};
  \item the Euclidean theory satisfies the standard transfer-matrix
    hypotheses used in constructive QFT: a reflection-invariant,
    hyperplane-factorizing Euclidean measure associated with $S_E$;
  \item Euclidean time evolution is implemented by the
    positivity-preserving semigroup of Core Axiom~\ref{ax:positivity}.
  \end{enumerate}
  Then for any $F \in \alg^{\mathrm{phys}}_+$ supported at Euclidean
  times $t > t_{\min} > 0$,
  \begin{equation}
    \langle \Theta F, F \rangle_E
    = \langle F | e^{-2t_{\min}H} | F\rangle_\hilb.
  \end{equation}
\end{proposition}
\begin{proof}
  Split the Euclidean field into negative-time, boundary, and
  positive-time parts:
  $\Phi = (\Phi_-, \Phi_0, \Phi_+)$. By
  Lemma~\ref{lem:hyperplane} and the factorization hypothesis in~(ii),
  the Euclidean measure can be written in the standard reflected form
  \[
    d\mu_E(\Phi)
    = Z^{-1} \,
      d\mu_-(\Phi_- \mid \Phi_0)\,
      e^{-B[\Phi_0]}\, d\nu(\Phi_0)\,
      d\mu_+(\Phi_+ \mid \Phi_0),
  \]
  where $\mu_-$ is the reflection of $\mu_+$.

  Because $F$ is supported in $x_0 > t_{\min}$, only the positive-time
  half-space enters $F$. Define the boundary-to-bulk map
  \[
    (K_{t_{\min}}F)(\Phi_0)
    := \int F(\Phi_+)\, d\mu_+^{\,t_{\min}}(\Phi_+ \mid \Phi_0),
  \]
  where $\mu_+^{\,t_{\min}}$ is the positive-time conditional measure
  truncated at distance $t_{\min}$ from the reflection hyperplane.
  Reflection invariance then gives
  \[
    \langle \Theta F, F\rangle_E
    = \int \overline{K_{t_{\min}}F(\Phi_0)}\,
      K_{t_{\min}}F(\Phi_0)\,
      e^{-B[\Phi_0]} d\nu(\Phi_0)
    = \|K_{t_{\min}}F\|_{L^2(\nu_B)}^2,
  \]
  where $d\nu_B := e^{-B[\Phi_0]} d\nu(\Phi_0)$.

  The standard Osterwalder--Schrader transfer-operator construction
  identifies $K_{t_{\min}}^*K_{t_{\min}}$ with the transfer operator
  $T_{2t_{\min}}$. Hypothesis~(iii) states that this transfer operator
  is implemented on the GNS Hilbert space by the positivity-preserving
  semigroup $e^{-tH}$. Therefore
  \[
    \langle \Theta F, F\rangle_E
    = \langle F, T_{2t_{\min}}F\rangle_\hilb
    = \langle F, e^{-2t_{\min}H}F\rangle_\hilb,
  \]
  which is the claimed identity.
\end{proof}

\begin{remark}[Status of OS2]
  The framework isolates the structural ingredients behind reflection
  positivity but does not claim that every local generating structure
  automatically yields a transfer matrix. The transfer-matrix
  hypotheses remain supplementary analytic input.
\end{remark}


\subsection{The Upstream Dynamical Chain for OS3}

The abstract framework isolates the exact clustering hypothesis needed
for OS3 and keeps model-dependent dynamical routes to that hypothesis
separate.

\begin{proposition}[Translation mixing implies OS3]\label{prop:mixing}
  Let $\pi$ be an $E(4)$-invariant stationary Euclidean measure in the
  flat sector, and let $\mathcal{B} \subset L^\infty(\pi)$ be a
  unital algebra of observables that is closed under translations:
  $\tau_a \mathcal{B} \subset \mathcal{B}$ for all
  $a \in \real^4$. Assume that the Schwinger functions are obtained as
  expectations of finite products of observables from~$\mathcal{B}$,
  and that for all $F, G \in \mathcal{B}$,
  \begin{equation}\label{eq:mixing}
    \big|\langle F \, \tau_a G\rangle_\pi
    - \langle F\rangle_\pi \langle G\rangle_\pi\big|
    \xrightarrow{|a| \to \infty} 0.
  \end{equation}
  Then the associated Schwinger functions satisfy OS3.
\end{proposition}
\begin{proof}
  Let
  \[
    F = O_{f_1}\cdots O_{f_m},
    \qquad
    G = O_{g_1}\cdots O_{g_n},
  \]
  with each factor in~$\mathcal{B}$. Since $\mathcal{B}$ is an algebra
  and is translation-stable, we have $F, G, \tau_a G \in \mathcal{B}$.
  By definition of the Schwinger functions,
  \[
    \langle F \, \tau_a G\rangle_\pi
    = S_{m+n}(f_1,\ldots,f_m,\tau_a g_1,\ldots,\tau_a g_n),
  \]
  while
  \[
    \langle F\rangle_\pi = S_m(f_1,\ldots,f_m),
    \qquad
    \langle G\rangle_\pi = S_n(g_1,\ldots,g_n).
  \]
  Applying~\eqref{eq:mixing} therefore gives
  \[
    S_{m+n}(f_1,\ldots,f_m,\tau_a g_1,\ldots,\tau_a g_n)
    \xrightarrow{|a|\to\infty}
    S_m(f_1,\ldots,f_m)\, S_n(g_1,\ldots,g_n),
  \]
  which is exactly the Osterwalder--Schrader clustering axiom OS3 for
  observables generated by~$\mathcal{B}$.
\end{proof}

\begin{remark}[A model-dependent spectral-gap route]
  In concrete theories, one standard route to
  \eqref{eq:mixing} is a spectral gap for an appropriate Euclidean
  transfer operator or Markov generator together with additional
  information identifying translated observables with the corresponding
  dynamical correlations. That route is useful, but it is not a
  theorem of the abstract framework and is therefore kept separate
  from Proposition~\ref{prop:mixing}.
\end{remark}

\begin{remark}[Axis~D refines the rate]
  OS3 asks only for factorization at large separation. Under D1, one
  expects exponential clustering once a quantitative form of
  \eqref{eq:mixing} is established. Under D2, one expects polynomial
  clustering. The
  \emph{existence} of OS3 is independent of Axis~D; the \emph{rate}
  depends on it.
\end{remark}


\subsection{Derivation of OS0--OS4}

\begin{theorem}[Flat stationary specialization --- detailed form]
  \label{thm:flat-detail}
  In the flat stationary sector, a system satisfying Core
  Axioms~\ref{ax:locality}--\ref{ax:positivity} together with branch
  choices on Axes A, B, and C produces Schwinger functions satisfying:
  \begin{enumerate}[nosep,label=(\roman*)]
  \item OS0 from Core Axiom~\ref{ax:resolution};
  \item OS1 from Core Axiom~\ref{ax:action} in the flat Euclidean
    sector, under the hypothesis that the equilibrium measure is
    uniquely $E(4)$-invariant;
  \item OS2 from Core Axioms~\ref{ax:action} and~\ref{ax:positivity},
    under the transfer-matrix hypotheses of
    Proposition~\ref{prop:transfer};
  \item OS4 from the chosen branch on Axis~B.
  \end{enumerate}
  If, in addition, the translation-mixing estimate
  \eqref{eq:mixing} holds for the observable algebra generating the
  Schwinger functions, then OS3 holds by
  Proposition~\ref{prop:mixing}. OS reconstruction then yields
  Wightman functions satisfying W0--W4.
\end{theorem}

\begin{proof}
  \textbf{OS0.} By Core Axiom~\ref{ax:resolution}, smeared
  observables satisfy~\eqref{eq:poly-bounds}. Therefore the Schwinger
  functions define continuous multilinear functionals on Schwartz test
  functions, hence tempered distributions on $(\real^4)^n$.

  \textbf{OS1.} In the flat sector, the local generating density
  depends on $\Phi$, $D\Phi$, and $\delta_{\mu\nu}$. Hence $S$ is
  $E(4)$-invariant, $e^{-S}$ is $E(4)$-invariant, and the equilibrium
  measure (uniquely selected by hypothesis) inherits that invariance.
  Expectations with respect to $\pi$ are therefore $E(4)$-covariant.

  \textbf{OS2.} For $F \in \alg^{\mathrm{phys}}_+$,
  Proposition~\ref{prop:transfer} gives
  $\langle \Theta F, F \rangle_E = \|e^{-t_{\min}H}F\|^2_\hilb \ge 0$,
  since $H \ge 0$ by Core Axiom~\ref{ax:positivity}.

  \textbf{OS4.} The grading is determined by Axis~B.

  \textbf{OS3.} Under the translation-mixing hypothesis,
  Proposition~\ref{prop:mixing} applies directly.

  \textbf{Wightman reconstruction.} Once OS0--OS4 are established, the
  OS reconstruction theorem~\cite{OS1, OS2} yields a Hilbert space,
  vacuum vector, operator-valued tempered fields, and a positive-energy
  representation of the Euclidean continuation, hence the Wightman
  axioms W0--W4. Core Axiom~\ref{ax:locality} identifies the upstream
  algebraic origin of locality, and Core Axiom~\ref{ax:propagation}
  explains why causal propagation is built into the Euclidean
  framework before reconstruction (Remark~\ref{rem:propagation-chain}).
\end{proof}


\subsection{General Manifold Specialization}

\begin{theorem}[Manifold-covariant specialization --- detailed form]
  \label{thm:manifold-detail}
  On a globally hyperbolic Lorentzian manifold, a system satisfying
  Core Axioms~\ref{ax:locality}--\ref{ax:positivity}, a choice on
  Axis~A, and geometric naturality in the sense of
  Remark~\ref{rem:naturality} induces a net of physical local
  $*$-algebras satisfying:
  \begin{enumerate}[nosep,label=(\alph*)]
  \item Isotony and locality, from Core Axioms~\ref{ax:locality}
    and~\ref{ax:propagation};
  \item Physical-algebra selection, from Axis~A;
  \item Local covariance, from Core Axiom~\ref{ax:action} together
    with Remark~\ref{rem:naturality}.
  \end{enumerate}
  If, in addition, the local algebra is generated by on-shell fields
  modulo the equations of motion and those equations are normally
  hyperbolic or Dirac-type, then the time-slice property holds. If the
  physical states are Hadamard, then the microlocal spectrum condition
  holds.
\end{theorem}

\begin{proof}
  Isotony is exactly clause~(i) of Core Axiom~\ref{ax:locality}, and
  locality is clause~(ii), with causal disjointness computed in the
  Lorentzian metric and controlled dynamically by Core
  Axiom~\ref{ax:propagation}. Axis~A specifies which local algebra is
  physical: the invariant subalgebra in the gauge branches and the
  full algebra in the ungauged branch.

  For local covariance, let $\psi: M_1 \hookrightarrow M_2$ be an
  orientation-preserving isometric embedding. By
  Remark~\ref{rem:naturality}, pullback of fields by $\psi$ induces an
  injective $*$-homomorphism
  \[
    \alpha_\psi: \alg_{M_1}(O) \to \alg_{M_2}(\psi(O))
  \]
  compatible with composition and identities. The assignment
  $M \mapsto \alg_M$ together with $\psi \mapsto \alpha_\psi$ is
  therefore a covariant functor on the spacetime category.

  If the local algebra is generated by on-shell fields modulo the
  equations of motion and those equations are normally hyperbolic or
  Dirac-type, then data on any neighborhood of a Cauchy surface
  determine the solution uniquely on its domain of dependence. Hence
  every observable in the larger region is represented by data already
  generated in the time-slice, so the time-slice property holds. If
  the physical states are Hadamard, then their two-point functions have
  the standard microlocal wavefront set, and the microlocal spectrum
  condition follows by the usual Radzikowski
  characterization~\cite{Radzikowski}.
\end{proof}


%% ===================================================================
\section{Compatibility, Non-Redundancy, and Separating Examples}
\label{sec:compatibility}
%% ===================================================================

\subsection{Compatibility}

\begin{remark}
  Compatibility is part of the architecture rather than a separate
  theorem: the core axioms constrain locality, resolution, propagation,
  generation, and positivity, while the specialization axes classify
  symmetry, statistics, interaction type, and infrared phase. No axis
  changes the meaning of a core axiom, and no core axiom fixes a
  branch value. The examples below record how the same core can be
  combined with different branch choices.
\end{remark}


\subsection{Non-Redundancy of the Specialization Axes}

The following are separating examples at the intended structural level;
they are not model-theoretic independence proofs and are stated as
examples rather than theorems for exactly that reason.

\begin{example}[Axis~A distinguishes gauge from ungauged theories]
  $\lambda\phi^4$ (Branch~A2) and pure $SU(N)$ Yang--Mills (Branch~A1)
  both satisfy the core axioms at the intended structural level but
  differ on Axis~A.
\end{example}

\begin{example}[Axis~B distinguishes statistics sectors]
  A free massive scalar (B1) and a free Dirac field (B2) satisfy the
  same core axioms and Branches A2, C2, but differ on Axis~B.
\end{example}

\begin{example}[Axis~C distinguishes free from interacting theories]
  The free massive scalar (C2) and $\lambda\phi^4$ (C1) satisfy the
  same core axioms and Branches A2, B1, but differ on Axis~C.
\end{example}

\begin{example}[Axis~D distinguishes infrared phases]
  A scalar theory on A2, B1, C1 with tunable parameters gives D1
  (gapped) at generic parameters and D2 (critical) at a tuned critical
  point, with the same core axioms and A/B/C branch data.
\end{example}


\subsection{Core Separating Examples}

\begin{remark}
  The following examples illustrate why no core axiom should be
  regarded as obviously redundant.
\end{remark}

\begin{example}[Locality can fail while the other structures remain]
  A theory with a single global observable algebra and no
  region-by-region assignment may still have finite resolution, finite
  propagation, additive generation, and positivity, but it has no
  isotony/locality net and therefore fails Core
  Axiom~\ref{ax:locality}.
\end{example}

\begin{example}[Finite resolution can fail while locality survives]
  A continuum theory with uncontrolled ultraviolet singularities may
  still have local algebras, finite propagation, local generation, and
  positive energy, but it fails the polynomial temperedness bounds of
  Core Axiom~\ref{ax:resolution}.
\end{example}

\begin{example}[Finite propagation can fail]
  Brownian motion has instantaneous spreading and therefore violates
  finite information speed while still serving as an example of local
  observables with additive generation and positivity.
\end{example}

\begin{example}[Additive decomposition can fail]
  A bilocal or mean-field action violates the separated-region
  additivity required in Core Axiom~\ref{ax:action} while leaving the
  other structural notions meaningful.
\end{example}

\begin{example}[Positivity can fail]
  A wrong-sign kinetic term yields an energy form unbounded below even
  if one still writes down local observables, finite resolution data,
  and local equations of motion.
\end{example}


%% ===================================================================
\section{Constructiveness, Renormalization, and Scaling}
\label{sec:constructive}
%% ===================================================================

\subsection{Why ``Constructive''}

The term is used in contrast with downstream axiomatizations. The
Wightman axioms ask for Hilbert-space data; the OS axioms ask for
Schwinger functions; the Haag--Kastler axioms ask for nets. In all
cases one verifies properties of an already-produced theory.

The present axioms ask for direct structural properties of the
generating dynamical system: does it admit a net of local algebras?
Does it have finite operational resolution? Does it propagate
information at finite speed? Is its generating structure local with
additive decomposition? Is the generator positive?

It is important to distinguish this usage from the standard meaning of
``constructive'' in mathematical physics (as in Glimm--Jaffe), which
refers to explicit constructions with detailed analytic estimates. The
present paper isolates the \emph{structural} layer of a constructive
program. The analytic layer --- where one proves existence and
properties of concrete theories using whatever techniques are
appropriate --- is logically separate from the framework itself.


\subsection{What the Axioms Prove and What They Do Not}

The axioms describe a constructive theory at one positive scale. A
continuum family $\{\mathcal{T}_\varepsilon\}_{\varepsilon > 0}$ with
$\ell_L(\varepsilon) \to 0$ requires an additional limiting theorem.

The core axioms do not by themselves prove that a continuum limit
exists, that a critical point exists, that a spectral gap exists, or
that transfer-matrix factorization holds. They sharply separate these
analytic tasks from the structural question of whether the underlying
dynamical system has the right architecture.


\subsection{Renormalization Group Within the Framework}

Let $\mathcal{T}_{\ell_L}$ denote a theory at scale $\ell_L$.
Coarse-graining to $\ell' > \ell_L$ gives
$\mathcal{T}_{\ell_L} \mapsto \mathcal{T}_{\ell'}$ with modified
couplings but the same structural type. Each member of this scale
family is evaluated against the same core axioms and branch logic.

If a theory is constructed by cutoff removal, the bare parameters may
depend on the cutoff. If it is constructed at fixed $\ell_L$ and
analyzed through a different limiting procedure, the physically
relevant running concerns effective couplings extracted after
coarse-graining, while the microscopic resolution stays positive.


%% ===================================================================
\section{Supersymmetry, Topological Theories, and the Boundary of Scope}
\label{sec:susy}
%% ===================================================================

\subsection{Supersymmetry Fits the Framework Cleanly}

Supersymmetry creates no friction with the five core axioms. The
correct structural place is Branch~B4. The key constructive gains are:

\begin{enumerate}[nosep]
\item \textbf{Positivity becomes internally certified.} Under B4,
  $H = Q^\dagger Q$ gives $H \ge 0$ automatically
  (Proposition~\ref{prop:susy-positivity}). This reduces the
  verification burden for Core Axiom~5 in the supersymmetric sector.
\item \textbf{Vacuum protection.} The Witten index is stable under
  continuous deformations~\cite{Witten}. When nonzero, the vacuum
  energy vanishes exactly.
\item \textbf{Transfer-matrix strengthening.} Once Euclidean time
  evolution is identified with $e^{-tH}$, the factorization
  $H = Q^\dagger Q$ gives an additional internal reason that the
  transfer semigroup is positive.
\end{enumerate}


\subsection{Pure TQFTs Mark a Genuine Scope Boundary}\label{sec:tqft}

Purely topological quantum field theories create real friction with
the present axioms. The core axioms were designed for QFTs with local
propagating degrees of freedom. Pure TQFTs do not have such degrees
of freedom. As a result, several core axioms become vacuous or
degenerate:

\begin{itemize}[nosep]
\item \textbf{Core Axiom~1.} A TQFT can assign finite-dimensional
  algebras to regions, but the dependence is topological rather than
  metric. The axiom holds only degenerately.
\item \textbf{Core Axiom~2.} Finite resolution is conceptually
  mismatched: polynomial bounds hold trivially because the spaces are
  finite-dimensional, but they control no nontrivial ultraviolet
  structure.
\item \textbf{Core Axiom~3.} Finite propagation is vacuous when there
  are no propagating local excitations.
\item \textbf{Core Axiom~4.} This is the least problematic: Chern--Simons
  and BF theories admit local action functionals.
\item \textbf{Core Axiom~5.} Often $H = 0$, so the axiom holds
  vacuously.
\end{itemize}

Pure TQFTs lie at the boundary of scope rather than in the main body.
For them, the Atiyah--Segal functorial formulation~\cite{Atiyah} is
the more natural axiomatic language.


\subsection{Topological Sectors Inside Dynamical Theories}

The framework does capture topological sectors of dynamical theories
whenever those sectors arise inside a theory with propagating local
degrees of freedom. Examples: $\theta$-vacuum structure in Yang--Mills,
topological protection of the Witten index. What lies outside scope
is not topological information as such, but theories that are purely
topological with no propagating sector at all.


%% ===================================================================
\section{Discussion and Outlook}
\label{sec:discussion}
%% ===================================================================

\subsection{What Is New}

The main advances of the present version are:

\begin{enumerate}[nosep]
\item The detailed correspondence between the five core axioms and
  each of the three standard axiomatic frameworks (Section~\ref{sec:existing}),
  with explicit identification of which core axiom produces which
  downstream property and where additional hypotheses are needed.

\item The reformulation of Core Axiom~4 as ``local dynamics with
  additive decomposition'' rather than ``local action functional,''
  making the framework broad enough to accommodate local generators
  beyond the classical Lagrangian case
  (Remark~\ref{rem:stochastic-action}).

\item The replacement of the overstrong OS3 claim by the exact
  theorem-level mixing hypothesis \eqref{eq:mixing} together with a
  full proof that this hypothesis yields OS3
  (Proposition~\ref{prop:mixing}), while keeping spectral-gap
  arguments in their proper model-dependent role.

\item The explanation of why Core Axiom~3 (finite propagation) is
  essential for the Euclidean-to-Lorentzian transition: it provides
  causal structure in the Euclidean branch, justifies the Wick
  rotation, and produces Lorentz invariance and the spectral condition
  in the reconstructed Lorentzian theory
  (Remark~\ref{rem:propagation-chain}).

\item Integration of supersymmetry as Branch~B4 and clear treatment of
  TQFTs as a scope boundary.

\item The explicit addition of geometric naturality
  (Remark~\ref{rem:naturality}) as the hypothesis that closes the gap
  between local generation and genuine locally covariant functoriality.

\item Positioning relative to the active constructive QFT programs of
  Hairer, Gubinelli--Imkeller--Perkowski, and the Balaban program.
\end{enumerate}


\subsection{What Is Not Claimed}

This paper does not claim that the five core axioms characterize all
quantum field theories, nor that the transfer-matrix, Hadamard,
spectral-gap, or critical-scaling hypotheses follow automatically from
them. It does not claim that every continuum limit exists, that every
critical point exists, or that every perturbative series can be
promoted to a full constructive theory.

The paper is a framework paper: it isolates upstream structural
conditions and shows how the standard axiomatic frameworks reappear
downstream once the corresponding additional analytic hypotheses are
imposed.


\subsection{Further Refinements}

\begin{description}[nosep]
\item[Gapless nonconformal spectral branch.] Theories with massless
  gauge bosons suggest a future Branch~D3.
\item[Extended supersymmetry.] Branch~B4 captures the basic
  supersymmetric strengthening. Extended SUSY, twisting, and other
  refinements can be layered on top.
\item[Topological sectors inside dynamical theories.] The right
  direction is not to force pure TQFTs into the core but to analyze
  topological sectors (e.g., $\theta$-vacua) inside propagating
  theories.
\item[Relation to singular SPDE methods.] The framework should be
  compatible with singular-SPDE methods such as Hairer's regularity
  structures~\cite{Hairer2014} or paracontrolled
  distributions~\cite{GIP}. Making that compatibility fully explicit
  would sharpen the connection.
\end{description}


%% ===================================================================
\section{Conclusion}
\label{sec:conclusion}
%% ===================================================================

The present framework proposes five core axioms shared by all
constructive quantum field theories with propagating local degrees of
freedom, together with four specialization axes. The core axioms
supply locality, finite resolution, finite propagation, local dynamics
with additive decomposition, and positivity. The axes classify
symmetry, statistics, interaction type, and infrared phase.

The framework's value lies not in proving new results about specific
models but in organizing the logical dependencies cleanly. In the flat
stationary sector, the core axioms plus branch data and
supplementary analytic hypotheses yield the OS axioms and hence the
Wightman axioms. In the Lorentzian manifold branch, they yield the
Haag--Kastler / locally covariant AQFT structure. The places where the
framework should not be stretched --- such as pure TQFTs or theories
without any local generating structure --- are stated openly.

The paper is therefore an abstract framework statement: it proposes an
organizational architecture, maps it in detail to the standard
frameworks, and separates the universal structural inputs from the
additional analytic hypotheses required for concrete applications.


%% ===================================================================
%% References
%% ===================================================================
\begin{thebibliography}{99}

\bibitem{Wightman1956}
A.~S. Wightman,
``Quantum field theory in terms of vacuum expectation values,''
\emph{Physical Review}, 101(2):860--866, 1956.

\bibitem{StreaterWightman}
R.~F. Streater and A.~S. Wightman,
\emph{PCT, Spin and Statistics, and All That},
W.~A. Benjamin, 1964.

\bibitem{OS1}
K. Osterwalder and R.~Schrader,
``Axioms for Euclidean Green's functions,''
\emph{Commun.\ Math.\ Phys.}, 31:83--112, 1973.

\bibitem{OS2}
K. Osterwalder and R.~Schrader,
``Axioms for Euclidean Green's functions.\ II,''
\emph{Commun.\ Math.\ Phys.}, 42:281--305, 1975.

\bibitem{GlimmJaffe}
J.~Glimm and A.~Jaffe,
\emph{Quantum Physics: A Functional Integral Point of View},
Springer, 2nd edition, 1987.

\bibitem{Haag}
R.~Haag,
\emph{Local Quantum Physics: Fields, Particles, Algebras},
Springer, 2nd edition, 1996.

\bibitem{BFV}
R.~Brunetti, K.~Fredenhagen, and R.~Verch,
``The generally covariant locality principle ---
a new paradigm for local quantum field theory,''
\emph{Commun.\ Math.\ Phys.}, 237:31--68, 2003.

\bibitem{Radzikowski}
M.~J. Radzikowski,
``Micro-local approach to the Hadamard condition in quantum field
theory on curved spacetime,''
\emph{Commun.\ Math.\ Phys.}, 179:529--553, 1996.

\bibitem{Leimkuhler}
B.~Leimkuhler and C.~Matthews,
\emph{Molecular Dynamics: With Deterministic and Stochastic Numerical
  Methods}, Springer, 2015.

\bibitem{Wilson}
K.~G. Wilson,
``Confinement of quarks,''
\emph{Physical Review D}, 10(8):2445--2459, 1974.

\bibitem{Simon}
B.~Simon,
\emph{Functional Integration and Quantum Physics},
AMS Chelsea Publishing, 2nd edition, 2005.

\bibitem{Witten}
E.~Witten,
``Constraints on supersymmetry breaking,''
\emph{Nuclear Physics B}, 202(2):253--316, 1982.

\bibitem{Atiyah}
M.~Atiyah,
``Topological quantum field theories,''
\emph{Publ.\ Math.\ IH\'ES}, 68:175--186, 1988.

\bibitem{Hairer2014}
M.~Hairer,
``A theory of regularity structures,''
\emph{Inventiones Mathematicae}, 198(2):269--504, 2014.

\bibitem{GIP}
M.~Gubinelli, P.~Imkeller, and N.~Perkowski,
``Paracontrolled distributions and singular PDEs,''
\emph{Forum of Mathematics, Pi}, 3:e6, 2015.

\bibitem{Balaban1}
T.~Balaban,
``Ultraviolet stability of three-dimensional lattice pure gauge field
theories,'' \emph{Commun.\ Math.\ Phys.}, 102(2):255--275, 1985.

\bibitem{Balaban2}
J.~Dimock and T.~R. Hurd,
``Construction of the two-dimensional Gross--Neveu model,''
\emph{Commun.\ Math.\ Phys.}, 137:263--287, 1991.

\bibitem{Villani2009}
C.~Villani,
``Hypocoercivity,''
\emph{Memoirs of the AMS}, 202(950), 2009.

\bibitem{FractalAI}
G.~Duran Ballester and S.~Hern\'andez Cerezo,
``Fractal AI: A fragile theory of intelligence,''
arXiv:1803.05049, 2018.

\end{thebibliography}

\end{document}
